\frontmatter% Use roman page numbering style (i, ii, iii, iv...) for the pre-content pages

\pagestyle{plain} % Default to the plain heading style until the thesis style is called for the body content

%----------------------------------------------------------------------------------------
%	TITLE PAGE
%----------------------------------------------------------------------------------------

\begin{titlepage}
\begin{center}
\begin{spacing}{1.2}

{\huge \bfseries \ttitle\par}

\vspace{5mm}
\textit{A thesis submitted in partial fulfillment of the requirements\\for the award of the degree of} 

\vspace{7mm}
\textsc{\huge \MakeUppercase{\degreename}}

\vspace{3mm}
\centering
\includegraphics[width=5cm]{Figures/IIITA_Logo.png}
\par\vspace{5mm}

\noindent\makebox[\textwidth][c]{%
\begin{minipage}[t]{0.45\textwidth}
    \begin{flushleft} \large
        \textit{By:-} \\
        \textsc{\studentnamesformatted}
    \end{flushleft}
\end{minipage}
\hfill
\begin{minipage}[t]{0.45\textwidth}
    \begin{flushright} \large
        \textit{Enrollment No.} \\
        \textsc{\rollnumbersformatted}
    \end{flushright}
\end{minipage}}\\[1cm]

\textit{Under the Supervision of}\\[2mm]
\textsc{\Large \supname}\\
\vspace{7mm}
\textit{to the}\\[2mm]
\textsc{\Large \deptname}\\

\vspace{8mm}
\textsc{\Large \univname}

\vspace{5mm}
November 2026

\end{spacing}
\end{center}
\end{titlepage}

%----------------------------------------------------------------------------------------
%	CERTIFICATE
%----------------------------------------------------------------------------------------
\checktoopen
\begin{figure}[htp]
    \includegraphics[height=3cm,keepaspectratio]{Figures/IIITA_Header.png}
\end{figure}
\thispagestyle{empty}
\vspace*{.04\textheight}

\begin{spacing}{1.5}
\begin{center}
    {\centering\large\bfseries CERTIFICATE \par\vspace{10pt}}
\end{center}

\noindent It is certified that the work contained in the thesis titled \textbf{\enquote{\ttitle}} by \textbf{\studentnamesdeclaration} has been carried out under my supervision and that this work has not been submitted elsewhere for a degree.

\vspace{3.5cm}
\hfill\begin{minipage}{7.5cm}
    \begin{spacing}{1.2}
        \par
        \rule{\textwidth}{0.2pt}\\
        {\supname} \par
        {\deptname} \par
        IIIT Allahabad \par
    \end{spacing}
\end{minipage}

\end{spacing}
\cleardoublepage

%----------------------------------------------------------------------------------------
%	DECLARATION
%----------------------------------------------------------------------------------------
\checktoopen
\begin{figure}[htp]
    \includegraphics[height=3cm,keepaspectratio]{Figures/IIITA_Header.png}
\end{figure}
\thispagestyle{empty}
\vspace{1mm}

\begin{center}
    {\large\bfseries CANDIDATE DECLARATION}
\end{center}

\begin{spacing}{1.2}

\noindent \pronounwe, \studentnamesdeclaration, hereby certify that the thesis titled \textbf{\enquote{\ttitle}} has been submitted by \pronounme{} in partial fulfillment of the requirements for the Degree of \degreename{} in the \deptname{} at the \univname.

\par
\noindent \pronounwe{} acknowledge that plagiarism includes: 
\begin{enumerate} 
    \item Reproducing another person's work (fully or partially) or ideas and presenting them as our own.
    \item Copying or paraphrasing another person's work without appropriate attribution.
    \item Engaging in literary theft by reproducing a unique literary construct without crediting its source.
\end{enumerate}
\par
\pronounwe{} affirm that due credit has been given to all original authors and sources through proper citations for words, ideas, diagrams, graphics, computer programs, experiments, results, and websites that are not \pronounour{} own contributions. Direct quotations have been appropriately marked, and sources have been properly referenced. Furthermore, \pronounwe{} declare that no portion of \pronounour{} work is plagiarized. In the event of a plagiarism claim, \pronounwe{} accept full responsibility for the contents of this thesis. \pronounwe{} acknowledge that our supervisor may not be able to verify the originality of every aspect of \pronounour{} work.
\end{spacing}

\vspace{1.5cm}
\noindent
\begin{minipage}[t]{0.4\textwidth}  
    \begin{flushleft} 
        Place: Prayagraj\\[2mm]
        Date: November 2026
    \end{flushleft}
\end{minipage}
\hfill
\begin{minipage}[t]{0.5\textwidth}
    \begin{flushright}         
        \textbf{\studentonename} (\studentoneroll)\\[4mm]
        \textbf{\studenttwoname} (\studenttworoll)
    \end{flushright}
\end{minipage}

\cleardoublepage

%----------------------------------------------------------------------------------------
%	ABSTRACT PAGE
%----------------------------------------------------------------------------------------

\begin{abstract}
\begin{spacing}{1.5}

\addchaptertocentry{\abstractname}

Indian Sign Language (ISL) is the primary native visual language used by more than 18 million individuals with hearing and speech impairments across India. Communication between deaf signers and hearing individuals remains severely constrained due to an acute shortage of certified interpreters (fewer than 1 interpreter per 18,000 deaf persons) and a lack of accessible technological translation bridges. Unlike American Sign Language (ASL), which relies primarily on single-handed finger-spelling, ISL is predominantly \textbf{bimanual (over 76\% two-handed gestures)}, characterized by intricate spatial orientations, high inter-class visual similarity, and subtle fingertip contacts.

This thesis presents a production-grade, accessible, bidirectional Indian Sign Language translation platform operating inside modern web browsers on commodity camera hardware without requiring sensor gloves or proprietary depth cameras. The forward Sign-to-Text pathway captures webcam video frames, estimates 21 three-dimensional skeletal landmarks per hand using Google MediaPipe Hands, constructs a normalized \textbf{156-dimensional scale- and translation-invariant bimanual feature representation}, and classifies \textbf{45 gesture classes} (A--Z, digits 1--9, and conversational signs) using an ensemble Random Forest classifier, supported by an MLP baseline. A temporal hysteresis state machine featuring a 1.5-second hold confirmation window and lockout intervals eliminates high-frequency prediction flicker and accidental character insertions.

The reverse pathway provides a multimodal translation engine: a 204-entry sign dictionary backed by 203 demonstration videos with speed control ($0.5\times$ to $2.0\times$) and camera practice challenger; browser-recognized speech-to-sign translation (English and Hindi); and an automated YouTube-to-Sign transcription workflow using \texttt{yt-dlp}, Deepgram Nova-2 speech AI, and timestamp-aligned video demonstration sequencing. The research is grounded in an empirical survey of 18 academic research papers and supported by an interactive gallery of 189 extracted research figures.

The system achieves 85.2\% held-out test accuracy, 84.8\% macro F1-score, 1.2 ms classifier inference, and 27.2 ms end-to-end processing latency (sustaining $\approx$36.8 FPS) on commodity hardware, providing a robust communication bridge for educational and assistive applications.

\end{spacing}
\end{abstract}

%----------------------------------------------------------------------------------------
%	ACKNOWLEDGEMENTS
%----------------------------------------------------------------------------------------

\begin{acknowledgements}
\begin{spacing}{1.5}
\addchaptertocentry{\acknowledgementname}

I express my deepest gratitude to my project supervisor, \textbf{Prof. Satish Kumar Singh}, Department of Information Technology, Indian Institute of Information Technology Allahabad, for his invaluable guidance, inspiring mentorship, and constructive criticism throughout the execution of this B.Tech project. His domain expertise and high standards have profoundly shaped the quality of this work.

I also extend my sincere appreciation to the faculty members and laboratory staff of the Department of Information Technology for providing computational resources and an academically stimulating research environment.

Special gratitude is extended to the Indian Deaf and Hard-of-Hearing community and accessibility advocates whose linguistic richness inspired this initiative, and to the open-source software community for the computational frameworks that made this project possible.

\end{spacing}
\end{acknowledgements}

%----------------------------------------------------------------------------------------
%	LIST OF CONTENTS/FIGURES/TABLES PAGES
%----------------------------------------------------------------------------------------

\tableofcontents

\listoffigures

\listoftables

%----------------------------------------------------------------------------------------
%	ABBREVIATIONS
%----------------------------------------------------------------------------------------

\begin{abbreviations}{ll}

\textbf{ANN} & \textbf{A}rtificial \textbf{N}eural \textbf{N}etwork\\
\textbf{ASL} & \textbf{A}merican \textbf{S}ign \textbf{L}anguage\\
\textbf{ASR} & \textbf{A}utomatic \textbf{S}peech \textbf{R}ecognition\\
\textbf{BoVW} & \textbf{B}ag-of-\textbf{V}isual-\textbf{W}ords\\
\textbf{BSL} & \textbf{B}ritish \textbf{S}ign \textbf{L}anguage\\
\textbf{CNN} & \textbf{C}onvolutional \textbf{N}eural \textbf{N}etwork\\
\textbf{DHH} & \textbf{D}eaf and \textbf{H}ard of \textbf{H}earing\\
\textbf{FPS} & \textbf{F}rames \textbf{P}er \textbf{S}econd\\
\textbf{HCI} & \textbf{H}uman--\textbf{C}omputer \textbf{I}nteraction\\
\textbf{HMM} & \textbf{H}idden \textbf{M}arkov \textbf{M}odel\\
\textbf{HUD} & \textbf{H}eads-\textbf{U}p \textbf{D}isplay\\
\textbf{IMU} & \textbf{I}nertial \textbf{M}easurement \textbf{U}nit\\
\textbf{ISL} & \textbf{I}ndian \textbf{S}ign \textbf{L}anguage\\
\textbf{KNN} & \textbf{K}-\textbf{N}earest \textbf{N}eighbors\\
\textbf{LSTM} & \textbf{L}ong \textbf{S}hort-\textbf{T}erm \textbf{M}emory\\
\textbf{MLP} & \textbf{M}ulti-\textbf{L}ayer \textbf{P}erceptron\\
\textbf{NLP} & \textbf{N}atural \textbf{L}anguage \textbf{P}rocessing\\
\textbf{NPS} & \textbf{N}et \textbf{P}romoter \textbf{S}core\\
\textbf{REST} & \textbf{RE}presentational \textbf{S}tate \textbf{T}ransfer\\
\textbf{RF} & \textbf{R}andom \textbf{F}orest\\
\textbf{SLR} & \textbf{S}ign \textbf{L}anguage \textbf{R}ecognition\\
\textbf{SPA} & \textbf{S}ingle-\textbf{P}age \textbf{A}pplication\\
\textbf{SURF} & \textbf{S}peeded-\textbf{U}p \textbf{R}obust \textbf{F}eatures\\
\textbf{SVM} & \textbf{S}upport \textbf{V}ector \textbf{M}achine\\
\textbf{TTS} & \textbf{T}ext-\textbf{t}o-\textbf{S}peech\\
\textbf{WER} & \textbf{W}ord \textbf{E}rror \textbf{R}ate\\

\end{abbreviations}

%----------------------------------------------------------------------------------------
%	SYMBOLS
%----------------------------------------------------------------------------------------

\begin{symbols}{lll}

$\mathbf{L}_{h,i}$ & 3D coordinate vector of joint $i$ on hand $h$ & $\mathbb{R}^3$ \\
$\hat{\mathbf{L}}_{h,i}$ & Wrist-centered, scale-normalized landmark coordinate & $\mathbb{R}^3$ \\
$\mathbf{V}$ & Full bimanual feature vector (156 dimensions) & $\mathbb{R}^{156}$ \\
$\tau_{\text{hold}}$ & Temporal hysteresis hold threshold (1.5 s) & \text{seconds} \\
$\tau_{\text{conf}}$ & Gesture activation confidence threshold (0.85) & \text{scalar} \\
$\hat{y}$ & Predicted sign class label & $\mathcal{C}$ \\

\end{symbols}

%----------------------------------------------------------------------------------------
%	DEDICATION
%----------------------------------------------------------------------------------------

\dedicatory{Dedicated to the Indian Deaf and Hard-of-Hearing community and to researchers advancing universal accessibility through technology.}
